\chapter*{Sommario}
Questa tesi affronta il problema dell'affidabilità delle comunicazioni Vehicle-to-Everything (V2X) in scenari di guida cooperativa, con particolare riferimento al platooning. In tali contesti, la qualità del canale radio e la variabilità delle condizioni di propagazione possono compromettere la trasmissione tempestiva delle informazioni di controllo, con impatti diretti su stabilità, comfort e sicurezza del convoglio. L'obiettivo del lavoro è progettare e valutare un approccio multi-canale che integri più tecnologie di comunicazione, selezionando dinamicamente il canale più affidabile in funzione delle condizioni osservate.

La parte sperimentale unisce la simulazione a eventi discreti dei framework OMNeT++, Veins e Plexe con la simulazione a tempo continuo di Sumo. In primo luogo, viene descritto il processo di preparazione e configurazione dell'ambiente di simulazione, includendo la riproducibilità degli esperimenti e la gestione delle dipendenze software. Successivamente, viene importata e adattata nel nuovo ambiente la logica di Packet Delivery Ratio (PDR) e di SafeSwitch, originariamente sviluppata nel progetto Segata2021Plexe, insieme agli scenari sperimentali di riferimento.

Un contributo centrale del lavoro consiste nella migrazione del modulo plexe\_hetnet dall'utilizzo di SimuLTE a Simu5G, al fine di supportare tecnologie cellulari di nuova generazione e abilitare valutazioni più aderenti agli standard recenti. Infine, vengono eseguite campagne di simulazione per diversi controller di platooning e parametri di traffico, e vengono generati grafici e metriche quantitative (ad esempio, PDR e indicatori legati alla dinamica del convoglio) per confrontare le prestazioni delle diverse configurazioni e valutare l'efficacia delle politiche di selezione del canale.

I risultati mostrano come l'adozione di un approccio multi-canale con meccanismi di stima della qualità del link e commutazione controllata possa migliorare la robustezza della comunicazione in condizioni degradate, permettendo ai veicoli del platoon di mantenere una minore distanza tra loro. La tesi si conclude con un'analisi critica dei limiti del modello e con possibili estensioni future, tra cui l'introduzione di ulteriori canali, strategie di decisione più sofisticate e una validazione su scenari più complessi.